\chapter{Fazit und Ausblick}\label{ch:fazit}

% [Cross-Ref Konsolidierte Doku, K-I.2 2026-05-15]
% Architekturentscheidungen F1-F15 + F-EXTRA-1..8 (Forschungsfragen): docs/bausteine/04_architekturentscheidungen_F1_F15.md
% Architektur-Master REV7.7 (finale Antwort auf F2-F4): docs/architektur/02_aktueller_master_REV7_7.md
% Termin-8-Stand (Bilanz): docs/termine_konsolidiert/08_termin_8.md
% Master-Plan Konsolidierungs-Sprint: docs/MASTERPLAN_KONSOLIDIERUNG_TERMINE.md
% NOTICE Architekt-Direktive II 2026-05-14 (Lizenz-Reframing): docs/sessions/20260514-4400-v31-final-adapters-and-p27.md

\section{Beantwortung der Forschungsfragen}

\subsection*{F1 (Komposition)}

Die elf Permutations-Achsen (Page, Node, Traversal, ValueHandle,
Concurrency, Allocator, Prefetch, Telemetry, ISA, Layout, Reclamation)
spannen einen Konfigurations-Raum von $\sim$5,5 Milliarden Permutationen
auf. Stand-der-Technik-Implementierungen lassen sich durch
algorithm\_profile-XMLs eindeutig rekonstruieren: 30~SOTA-Profile
(8~Rang-1 + 22~Rang-2/3) belegen die Tragfaehigkeit der Achsen-
Konzeption fuer P01--P32.

Die hybride API-Spezialisierung (1~Param $\Rightarrow$
\texttt{std::vector}, 2~Params $\Rightarrow$ \texttt{std::map},
$N>2 \Rightarrow$ \texttt{std::map<K, std::tuple>}) erlaubt die
nahtlose Integration sowohl algorithmus-orientierter
(\texttt{std::map::find}/\texttt{insert}/\texttt{erase}) als auch
container-orientierter (\texttt{std::vector::push\_back}/\texttt{at})
Such-Patterns auf der gleichen Engine-Klasse.

\subsection*{F2 (Mess-Methodik)}

Die Mess-Methodik gliedert sich in zwei Modi:

\begin{itemize}
    \item \textbf{Production-Mode} (\texttt{COMDARE\_EXPERIMENT\_MODE=OFF},
          Default): kein Mess-Overhead, ResultAggregator nicht in der
          ExecutionEngine.
    \item \textbf{Experiment-Mode} (\texttt{COMDARE\_EXPERIMENT\_MODE=ON},
          vom messung\_driver gesetzt): ResultAggregator als
          integraler ExecutionEngine-Member, sammelt pro
          run\_workload-Aufruf die Hardware-Counter.
\end{itemize}

Profile-aware Workload-Routing (V11.2) waehlt pro Permutations-Modul
den Workload-Typ aus dem \texttt{traversal}-Tag des Profils, was
faire Vergleiche zwischen Range-Query- und Point-Lookup-orientierten
Algorithmen sicherstellt.

\subsection*{F3 (Pruefling-Vergleich)}

Der konkrete Vergleich von PRT-ART gegen die 30~SOTA-Profile steht
nach den ersten messung\_driver-Laeufen aus (siehe Ausblick~\S\ref{sec:ausblick}).
Die Architektur ist vollstaendig vorbereitet:

\begin{itemize}
    \item PRT-ART-Profil in
          \texttt{prt\_art/algorithm\_profiles/prtart\_pruefling.profile.xml}
    \item Drei Pflicht-Messreihen-Templates in
          \texttt{Code/experiment\_config/messreihen.xml}
    \item ExperimentDriver mit Auto-Pickup von 30~SOTA-Profilen
          (V10.5)
    \item PrtArtSearchEngineAdapter mit vollstaendiger ABI-Konformitaet
          inkl. notify\_*-Hooks (V11.1) und allen std::map-Operationen
          (V12.2 + V13.5 merge/extract)
\end{itemize}

\section{Limitierungen}

\begin{itemize}
    \item \textbf{Module-Bodies derzeit Stub:} Die Codegen-Pipeline
          generiert ABI-konforme Skelett-Module pro Profil. Die
          eigentliche Such-Algorithmus-Implementation pro Permutation
          bleibt fuer eine Folge-Phase (V15+) offen.
    \item \textbf{Hardware-Validierung ausstehend:} cmake configure ist
          gruen auf MSVC i7-1270P, aber der vollstaendige ctest-Lauf +
          messung\_driver-E2E auf BlockAO-Plattformen (X86\_AVX2/AVX512,
          ARM\_Neon/SVE) steht aus.
    \item \textbf{Manuskript-Auswertungs-Kapitel:} Aktuell methodisch
          beschrieben; konkrete Mess-Resultate folgen.
\end{itemize}

\section{Zukuenftige Arbeit}\label{sec:ausblick}

\textbf{Kurzfristig (V15+):}

\begin{itemize}
    \item E2E-Lauf der drei Pflicht-Messreihen auf realer Hardware
          (V15.2 in der Sprint-Plannung)
    \item Generierung der Mess-Diagramme via diagram\_generator und
          Integration in Kapitel~\ref{ch:auswertung}
    \item Vervollstaendigung der Module-Bodies pro Permutation
\end{itemize}

\textbf{Mittel- bis langfristig:}

\begin{itemize}
    \item Erweiterung um GPU-Adapter (Grace Hopper Cluster, ZIH TUD)
    \item PIM-Allokator-Integration (A16 Lee et al. HPCA 2026)
    \item Formale Verifikation einzelner Permutations-Klassen
          (StarMalloc-Style A13)
\end{itemize}

Die Drei-Repo-Architektur (cache-engine + prt-art + Diplomarbeit-Code)
ermoeglicht es zukuenftigen Diplomanden, weitere Pruefling-Algorithmen
analog zu PRT-ART einzubringen, ohne die Werkzeug-Bibliothek
(cache-engine) zu modifizieren --- ein wichtiges Designprinzip fuer
die Skalierbarkeit der Forschungs-Infrastruktur.

\section{Konsolidierungs-Beitraege (N + M-Phase 2026-05-18)}\label{sec:konsolidierungs-beitraege}

% Quelle: docs/sessions/20260518-4900-...md + 10_schichten_modell_M.md + 11_axes_vs_strategies_disambiguation.md + 07_bausteine_matrix_N_erweitert.md

Nach Termin~8 wurden zwei strukturelle Konsolidierungen vorgenommen,
die ueber den ursprueglichen Diplomarbeit-Scope hinausgehen:

\subsection{M-Modell: 4-Subsystem-Trennung}

Die Trennung von CacheEngineBuilder (autonomes
Plattform-Ausmess-System), CacheEngine (Werkzeug-Bibliothek) und
Pruefling (PRT-ART) als drei unabhaengige Subsysteme im selben
Repo, plus dem messung\_driver als vierter Outer-Loop, wurde
formalisiert (siehe Kapitel~\ref{sec:vier-subsysteme}). Die
\emph{bidirektionale} Cache-Engine~$\leftrightarrow$~Pruefling-Beziehung
(CE bindet Pruefling als Permutations-Struktur ein, Pruefling nutzt
CE als Optimierungs-Service) ist die zentrale Eigenschaft des
M-Modells. Diese Symmetrie war in REV7 implizit und wurde durch das
M-Modell explizit gemacht. Sie ermoeglicht die Multi-Pruefling-Faehigkeit
(Messreihe~A vergleicht PRT-ART gegen 22 V31-SOTA-Adapter parallel).

\subsection{N-Phase: 14-Achsen-Bausteine-Matrix}

Die Bausteine-Matrix wurde von 11~Hauptachsen auf
\textbf{14~Hauptachsen + ca.~30 Sub-Achsen} erweitert
(siehe Kapitel~\ref{sec:14-achsen-substruktur}). Die fuenf Achsen-
Aenderungen (3-Split, 6-Split, 8-Split, 11-Erweiterung, 12+13~NEU)
adressieren konzeptionelle Luecken der ursprueglichen Matrix:

\begin{itemize}
    \item Die fundamentale Trennung zwischen Algorithmus-Traversal und
    Cache-Memory-Traversal (Achse~3-Split) ist die Voraussetzung fuer
    eine saubere Mapping-Strategie zwischen den nicht-Cache-kohaerenten
    Konzepten des Entwicklers und der Cache-Mechanik der CacheEngine.

    \item Die feingliedrige Aufteilung der Allokations-Strategien
    (Achse~6.1--6.5) loest die vorher vermischten 7~Allokator-Achsen
    AA1--AA7 in 5~konzeptionell klar getrennte Sub-Achsen auf.

    \item Die NEUE Achse~13 SCHEDULING beruecksichtigt explizit die
    Hardware-Limitierung der SIMD-Einheiten (typisch 2~von $N$~Cores)
    und ermoeglicht damit \emph{heterogene} Worker-Pools, wo
    SIMD-faehige Worker und scalare Worker parallel laufen.
\end{itemize}

\subsection{PRT-ART als Bausteine-Spiegel}

Die N-Phase-Pflicht (User-Direktive 2026-05-18: "PRT-ART muss alle
Achsen detailliert benennen und konstruieren") qualifiziert PRT-ART
als \emph{neuartigen Algorithmus}: PRT-ART hat in den Achsen~1, 5,
6.5, 7, 8.1 und 10 jeweils Neu-Implementationen (die PRT-ART-spezifische
Bucket-Strategy mit 4+2 Pool-Multiplicity, die Distance-Estimator-
basierte Prefetch-Strategie, die OLC + Reserved-Value-Blocks-
Concurrency-Mechanik und die H1/H2/H3-Metriken). In den uebrigen
Achsen werden bestehende Stand-der-Technik-Bausteine via
Permutations-Konfiguration wiederverwendet.

\subsection{Anti-Vermischung der 4 Konzept-Ebenen}

Die Klarstellung der 4~unabhaengigen Konzept-Ebenen
(Ebene~I Bausteine-Achsen, Ebene~II CE-Sub-Engines, Ebene~III
Cache-Strategien, Ebene~IV Such-Engine-Familien) in
docs/architektur/11\_axes\_vs\_strategies\_disambiguation.md ist eine
methodische Konsequenz aus der Tieflektuere der 33~Quell-Paper:
Begriffe wie "Hardware-Probing", "Cache-Strategie" oder "Permutations-
Dimension" haben in der Literatur unterschiedliche Bedeutungen, was
die Ableitung der eigenen Bausteine-Matrix verschleiert hat. Die
4-Ebenen-Hierarchie ist die Basis fuer weitere Konsolidierungen.
